\documentclass[11pt]{article}
\usepackage[margin=1in]{geometry}
\usepackage{amsmath,amssymb,amsthm}
\usepackage{mathtools}
\usepackage{hyperref}

\newtheorem{theorem}{Theorem}
\newtheorem{lemma}[theorem]{Lemma}
\newtheorem{proposition}[theorem]{Proposition}
\theoremstyle{definition}
\newtheorem{definition}[theorem]{Definition}
\newtheorem{example}[theorem]{Example}

\newcommand{\inner}[2]{\langle #1, #2 \rangle}
\newcommand{\norm}[1]{\left\| #1 \right\|}
\newcommand{\R}{\mathbb{R}}
\newcommand{\C}{\mathbb{C}}
\newcommand{\F}{\mathbb{F}}

\title{Inner Product Spaces}
\author{Math Foundations --- Node 08}
\date{}

\begin{document}
\maketitle

\subsection*{First, an example}

Before any axioms, look at the standard dot product on $\R^2$. Take
$u = (3, 4)$ and $v = (-4, 3)$. Then
$$\inner{u}{v} = 3 \cdot (-4) + 4 \cdot 3 = -12 + 12 = 0,$$
\emph{Read aloud:} ``inner product of $u$ and $v$ equals three times negative four plus four times three, which is zero.''
so $u$ and $v$ are \emph{orthogonal}. Their lengths are
$$\norm{u} = \sqrt{3^2 + 4^2} = 5, \qquad \norm{v} = \sqrt{(-4)^2 + 3^2} = 5.$$
\emph{Read aloud:} ``norm of $u$ equals the square root of nine plus sixteen, which is five; likewise norm of $v$ equals five.''
So $u$ and $v$ are two perpendicular vectors of the same length $5$.
The whole chapter is about which abstract structures preserve this
``angle and length'' picture.

\section{Motivation}

In $\R^n$ the dot product $u \cdot v = \sum_i u_i v_i$ encodes both
length, via $\norm{u}^2 = u \cdot u$, and angle, via the relation
$\cos\theta = (u \cdot v)/(\norm{u}\,\norm{v})$. An \emph{inner product
space} is the abstraction that lets us carry these geometric notions
into any vector space, including spaces of functions, sequences, and
random variables. Linear algebra without an inner product gives us
linear maps; with one, it gives us geometry.

\section{Definitions}

Let $V$ be a vector space over $\F \in \{\R, \C\}$.

\begin{definition}[Inner product]
An \emph{inner product} on $V$ is a map
$\inner{\cdot}{\cdot} : V \times V \to \F$ satisfying, for all
$u, v, w \in V$ and $\alpha \in \F$:
\begin{enumerate}
  \item \textbf{Conjugate symmetry:} $\inner{u}{v} = \overline{\inner{v}{u}}$.
  \item \textbf{Linearity in the first argument:}
        $\inner{\alpha u + w}{v} = \alpha \inner{u}{v} + \inner{w}{v}$.
  \item \textbf{Positive definiteness:} $\inner{u}{u} \ge 0$, with
        equality iff $u = 0$.
\end{enumerate}
A vector space equipped with an inner product is called an
\emph{inner product space}.
\end{definition}

Note that over $\R$, conjugate symmetry collapses to symmetry, and
linearity in the first argument together with symmetry gives linearity
in the second argument as well. Over $\C$, the second argument is
\emph{conjugate}-linear; this is sometimes called the ``physicist's
convention'' when the conjugate sits in the first slot --- we follow
the mathematician's convention here.

\begin{definition}[Induced norm]
The \emph{induced norm} on an inner product space $(V, \inner{\cdot}{\cdot})$ is
$$\norm{v} := \sqrt{\inner{v}{v}}, \qquad v \in V.$$
\emph{Read aloud:} ``norm of $v$ is defined as the square root of the inner product of $v$ with itself, for every $v$ in $V$.''
Positive definiteness makes this well-defined and non-negative.
\end{definition}

\begin{definition}[Orthonormal set]
A finite or countable family $\{e_i\}_{i \in I} \subset V$ is called
\emph{orthonormal} if
$$\inner{e_i}{e_j} = \delta_{ij}
  = \begin{cases} 1 & i = j, \\ 0 & i \neq j.\end{cases}$$
\emph{Read aloud:} ``inner product of $e_i$ and $e_j$ equals the Kronecker delta $\delta_{ij}$, which is one when $i$ equals $j$ and zero otherwise.''
\end{definition}

\section{The Cauchy-Schwarz inequality}

\begin{theorem}[Cauchy-Schwarz]
\label{thm:cs}
Let $(V, \inner{\cdot}{\cdot})$ be an inner product space over $\F$.
For all $u, v \in V$,
\begin{equation}
\label{eq:cs}
\bigl| \inner{u}{v} \bigr| \;\le\; \norm{u}\,\norm{v},
\end{equation}
\emph{Read aloud:} ``the absolute value of the inner product of $u$ and $v$ is less than or equal to the norm of $u$ times the norm of $v$.''
with equality if and only if $u$ and $v$ are linearly dependent.
\end{theorem}

\begin{proof}
If $v = 0$ both sides of \eqref{eq:cs} are zero and there is nothing
to prove, so assume $v \neq 0$. The trick is to consider, for a scalar
$t \in \F$ to be chosen, the non-negative quantity
$$0 \le \norm{u - t v}^2 = \inner{u - t v}{u - t v}.$$
\emph{Read aloud:} ``zero is less than or equal to the squared norm of $u$ minus $t$ times $v$, which equals the inner product of that vector with itself.''
Expanding using linearity in the first slot and conjugate linearity in
the second,
\begin{align*}
\inner{u - tv}{u - tv}
  &= \inner{u}{u} - t\,\inner{v}{u} - \bar t\,\inner{u}{v} + |t|^2 \inner{v}{v} \\
  &= \norm{u}^2 - t\,\overline{\inner{u}{v}} - \bar t\,\inner{u}{v}
     + |t|^2 \norm{v}^2.
\end{align*}
\emph{Read aloud:} ``the inner product of $u$ minus $tv$ with itself expands as the inner product of $u$ with $u$, minus $t$ times the inner product of $v$ with $u$, minus $t$-bar times the inner product of $u$ with $v$, plus the modulus of $t$ squared times the inner product of $v$ with $v$; equivalently, norm-$u$ squared minus $t$ times the conjugate of the inner product of $u$ and $v$, minus $t$-bar times the inner product of $u$ and $v$, plus the modulus of $t$ squared times norm-$v$ squared.''
Choose $t = \inner{u}{v} / \norm{v}^2$. Then $\bar t \inner{u}{v} =
|\inner{u}{v}|^2 / \norm{v}^2$ and $t\,\overline{\inner{u}{v}} =
|\inner{u}{v}|^2 / \norm{v}^2$, while $|t|^2 \norm{v}^2 =
|\inner{u}{v}|^2 / \norm{v}^2$. Substituting,
$$0 \le \norm{u}^2 - \frac{|\inner{u}{v}|^2}{\norm{v}^2}.$$
\emph{Read aloud:} ``zero is less than or equal to norm-$u$ squared minus the modulus squared of the inner product of $u$ and $v$, divided by norm-$v$ squared.''
Multiplying through by $\norm{v}^2 > 0$ and rearranging,
$$|\inner{u}{v}|^2 \le \norm{u}^2 \norm{v}^2,$$
\emph{Read aloud:} ``the modulus squared of the inner product of $u$ and $v$ is less than or equal to norm-$u$ squared times norm-$v$ squared.''
and taking square roots gives \eqref{eq:cs}.

For the equality case: equality forces $\norm{u - t v}^2 = 0$, so by
positive definiteness $u = tv$, i.e.\ $u$ and $v$ are linearly
dependent. Conversely if $u = \lambda v$ then both sides of
\eqref{eq:cs} equal $|\lambda|\,\norm{v}^2$.
\end{proof}

\begin{proposition}[Triangle inequality]
For all $u, v \in V$, $\norm{u + v} \le \norm{u} + \norm{v}$.
\end{proposition}

\begin{proof}
Expand $\norm{u+v}^2 = \norm{u}^2 + 2\,\Re\,\inner{u}{v} + \norm{v}^2$,
bound $\Re\,\inner{u}{v} \le |\inner{u}{v}| \le \norm{u}\,\norm{v}$ by
Theorem~\ref{thm:cs}, and take square roots.
\end{proof}

\section{Examples}

\begin{example}[Euclidean space]
$V = \R^n$, $\inner{u}{v} = \sum_{i=1}^n u_i v_i$. The induced norm is
the usual Euclidean length, and Cauchy-Schwarz is the classical
$\bigl(\sum u_i v_i\bigr)^2 \le \bigl(\sum u_i^2\bigr)\bigl(\sum v_i^2\bigr)$.
\end{example}

\begin{example}[$L^2$ on the unit interval]
$V = C[0,1]$ with $\inner{f}{g} = \int_0^1 f(x)\,g(x)\,dx$. Cauchy-Schwarz
becomes
$$\left|\int_0^1 f g\,dx\right|
   \le \left(\int_0^1 f^2\,dx\right)^{1/2}
       \left(\int_0^1 g^2\,dx\right)^{1/2},$$
\emph{Read aloud:} ``the absolute value of the integral from zero to one of $f$ times $g$ $dx$ is less than or equal to the square root of the integral of $f$ squared, times the square root of the integral of $g$ squared.''
a special case of H\"older's inequality with $p = q = 2$.
\end{example}

\section{Where this leads}

Cauchy-Schwarz is the inequality that lets us define $\cos\theta =
\inner{u}{v}/(\norm{u}\norm{v}) \in [-1, 1]$. The construction of
orthonormal bases via Gram-Schmidt, and orthogonal projection onto a
closed subspace, both rest on this single inequality. In ML, every
similarity score that takes the form of a dot product --- attention
logits, contrastive loss similarities, cosine similarity --- is a
direct corollary.

\end{document}
